%% kapa: exam=1 round=35 year=2024 period=1 form=A subject=real_estate subject_ko=부동산학원론 seq=31 number=111
%% source: Q-Net 감정평가사 기출문제 (cst00309) artlSeq=5213229 file=raw/1008/5213229/1교시 문제지 원본.pdf
%% conversion: kordoc 4.24.51 -> tools/gen1.py (auto); review=proofread-sonnet; flags=none
\begin{question}{111}
부동산투자의 수익과 위험에 관한 설명으로 옳지 않은 것은?
\begin{choices}
  \item 다양한 자산들로 분산된 포트폴리오는 체계적 위험을 감소시킨다.
  \item 위험회피형 투자자는 위험 증가에 따른 보상으로 높은 기대수익률을 요구한다.
  \item 동일한 자산들로 구성된 포트폴리오라도 자산들의 구성비중에 따라 포트폴리오의 수익과 위험이 달라진다.
  \item 시장상황에 대한 자산가격의 민감도가 높을수록 수익률의 표준편차는 커진다.
  \item 지분투자수익률은 지분투자자의 투자성과를 나타낸다.
\end{choices}
\end{question}
\answer{1}
